\documentclass[3p,11pt]{elsarticle}
\usepackage{amsmath,amssymb,amsthm}
\usepackage{booktabs}
\usepackage{lineno}
\usepackage[hidelinks]{hyperref}
\modulolinenumbers[5]
\journal{Operations Research Letters}

\newtheorem{theorem}{Theorem}
\newtheorem{lemma}[theorem]{Lemma}
\newtheorem{proposition}[theorem]{Proposition}
\newtheorem{corollary}[theorem]{Corollary}
\theoremstyle{remark}
\newtheorem{remark}[theorem]{Remark}
\theoremstyle{definition}
\newtheorem{definition}[theorem]{Definition}

\newcommand{\R}{\mathbb{R}}
\newcommand{\Z}{\mathbb{Z}}
\newcommand{\supp}{\operatorname{supp}}
\newcommand{\OPT}{\mathrm{OPT}}

\begin{document}
\begin{frontmatter}
\title{Fixed charges of arbitrary sign: what survives and what fails}
\author{Duy T. Nguyen\corref{cor1}}
\ead{duynguyen.das@gmail.com}
\cortext[cor1]{Corresponding author. ORCID: \href{https://orcid.org/0009-0004-0732-2093}{0009-0004-0732-2093}}
\address{Independent Researcher, Ho Chi Minh City, Vietnam}

\begin{abstract}
For integer activities, conditioning on the support makes a fixed-charge objective
affine. If every support-conditioned cell is integral, an optimal solution is a vertex
of its cell for arbitrary fixed charges and marginal rates. A totally unimodular
constraint matrix with integral data guarantees this condition. This
surviving property is weaker than the classical conclusions. A negative fixed charge
rewards opening an activity at its minimum positive level, which the vertex property does
not constrain. We give unique-optimum instances that violate the fractional-period
property in capacitated lot sizing and the positive-support forest properties in
fixed-charge transportation and flow. A four-period lot-sizing instance further has
unrestricted optimum $-8$, while every schedule in the subplan state space used by a
classical $O(T^4)$ recursion costs at least $-5$. Finally, the one-dimensional
charge-if-positive function is piecewise concave in Zangwill's sense exactly for a
non-negative fixed charge. Thus integrality preserves a weak vertex statement across the
sign change, while the stronger structures and an algorithmic state space built on them
can fail.
\end{abstract}

\begin{keyword}
fixed-charge problems \sep negative prices \sep total unimodularity \sep lot sizing
\sep integral polyhedra
\end{keyword}
\end{frontmatter}
\linenumbers

\section{Introduction}

What remains of classical fixed-charge structure when a charge becomes a payment? The
question is too broad without a precise cost model. We study the discrete
charge-if-positive function
\begin{equation}\label{eq:fc}
\varphi_j(0)=0,\qquad \varphi_j(x)=s_j+c_jx\quad(x>0),
\end{equation}
where $s_j$ is a fixed charge and $c_j$ is a marginal rate. Both can inherit the sign of
one price: $s_j=p_j\sigma_j$ and $c_j=p_j\gamma_j$, with
$\sigma_j,\gamma_j>0$. Wholesale electricity provides the motivating case. Its price can
fall below zero, and these events grow with renewable capacity \cite{seel21}. The fixed
charge is then received rather than paid.

The sign of $s_j$ decides the curvature. For $x>0$ and $0<\lambda<1$,
\begin{equation}\label{eq:chord}
\varphi_j(\lambda x)-\lambda\varphi_j(x)=s_j(1-\lambda),
\end{equation}
so $\varphi_j$ is concave on $[0,u_j]$ exactly when $s_j\ge0$. When $s_j<0$ it lies
strictly below its own chord, which makes it \emph{convex}; what \eqref{eq:fc} loses at
the origin is not convexity but \emph{lower semicontinuity}, since
$\liminf_{x\downarrow0}\varphi_j(x)=s_j<0=\varphi_j(0)$.

This sign change removes two hypotheses used by the structure theory: concavity and lower
semicontinuity. Florian and Klein \cite{fk71} assume concavity;
Love \cite{love73} and Swoveland \cite{swoveland75} assume piecewise concavity. Koca,
Yaman and Akt\"urk \cite{kya14} assume piecewise concavity together with lower
semicontinuity. Akbalik, Penz and Rapine \cite{apr15} invoke Love's property explicitly
under concave costs. The same model has a polyhedral literature \cite{ak05}, which does
not use the structural property at all; and the concave case with arbitrary capacities is
NP-hard \cite{flr80}, so the polynomial results are all for restricted capacities. Ou \cite{ou17} imposes no sign restriction in his model statement,
but the structural property he uses is attributed to Swoveland and is proved there for
piecewise concave costs, so the hypothesis is inherited by the proof.

Two qualifications prevent an overclaim. First, \cite{kya14} already give a repair:
their \S3.4 models minimum order quantities by setting $p_t(x)=\infty$ on
$(0,L)\cup(C,\infty)$, and solves that case in $O(n^6)$. With $L=1$ the effective domain
becomes $\{0\}\cup[1,C]$, which restores lower semicontinuity and leaves a piecewise
concave cost whose breakpoints $\{0,1,C\}$ are time-invariant and fixed in number. Their
polynomial algorithm requires precisely this condition. Negative charges are
therefore not outside their framework; they are outside the framework as it is usually
instantiated. The polynomial solvability of the $L=1$ case in \cite{kya14} is not at issue
here. What fails, and what Sections~\ref{sec:fails}--\ref{sec:alg} are about, is the
transfer of the classical structural properties, and of the algorithms built on them, to
charges of either sign.
Second, \cite{love73} is closer than it first appears. His model carries production bounds
$x^L_i\le x_i\le x^U_i$, and his conclusion excepts production equal to ``zero or its
upper or lower bound''. Setting $x^L_i=1$ gives the exception set $\{0,1,x^U_i\}$, exactly
the conclusion of Corollary~\ref{cor:count} below. What the corollary adds
is that this holds \emph{without} concavity. Love assumes his cost concave on
$[0,\infty)$, which \eqref{eq:fc} is not when $s_j<0$. The classical property in the usual
$x^L_i=0$ setting is a different statement, and Section~\ref{sec:fails} shows it fails.

The surviving statement has a short proof. Conditioning on the support leaves an
\emph{affine} function on a bounded integral polyhedron. We claim no depth for this step.
The point is that the classical route is closed, while this weaker route remains open.
Proposition~\ref{prop:zangwill} makes the boundary exact.
Zangwill \cite{zangwill67} is the canonical reference here, and it is how \cite{love73}
reaches the lot-sizing result. His theory is not confined to continuous costs: his
Theorem~6 minimises any \emph{piecewise concave} function over its dominant set, and such
functions may jump. Proposition~\ref{prop:zangwill} shows that \eqref{eq:fc} lies in that
class precisely when $s_j\ge0$, so the sign of the charge is the exact boundary of the
theory rather than an obstruction we happen to meet.

\paragraph{Contribution.} Four claims answer the opening question. First,
Theorem~\ref{thm:main} and Corollary~\ref{cor:count} identify the weak vertex statement
that survives arbitrary signs. Second, Section~\ref{sec:fails} gives three
unique-optimum witnesses that violate the stronger classical statements. Third,
Section~\ref{sec:alg} shows that the $O(T^4)$ state space of \cite{apr15}, correct under
its assumptions, excludes the unique optimum of a four-period signed-charge instance.
Finally, Proposition~\ref{prop:sharp} shows why support-cell integrality is needed for our
proof, and Remark~\ref{rem:model} records the associated modelling trap.

\section{The problem class}

Let $A$ be a real matrix, let $b$ be a real right-hand side, let
$u\in\Z_{>0}^n$, and consider
\begin{equation}\label{eq:P}
(\mathrm{FC})\qquad
\min\ \Phi(x)=\sum_{j=1}^{n}\varphi_j(x_j)
\quad\text{over}\quad x\in\Z^n,\quad Ax\lesseqgtr b,\quad 0\le x\le u,
\end{equation}
with $\varphi_j$ as in \eqref{eq:fc} and $s_j,c_j\in\R$ \emph{of arbitrary sign}.
Assume throughout that \eqref{eq:P} is feasible.
Write $\supp(x)=\{j:x_j>0\}$, and for $S\subseteq\{1,\dots,n\}$ define the \emph{cell}
\begin{equation}\label{eq:cell}
P(S)=\bigl\{x\in\R^n:\ Ax\lesseqgtr b,\quad x_j=0\ (j\notin S),\quad 1\le x_j\le u_j\ (j\in S)\bigr\}.
\end{equation}
The constant $1$ in \eqref{eq:cell} is where integrality enters, and it is the source of
everything in Section~\ref{sec:fails}: for continuous variables the set
$\{x_j>0,\ j\in S\}$ is not closed and carries no vertex guarantee.

\section{The vertex theorem}

\begin{lemma}\label{lem:cell}
Let $x^{*}$ be feasible for \eqref{eq:P} and $S=\supp(x^{*})$. Then $x^{*}\in P(S)$;
every integral point of $P(S)$ is feasible for \eqref{eq:P} with support exactly $S$;
$\Phi$ agrees on $P(S)$ with the affine function $\sum_{j\in S}s_j+\sum_{j\in S}c_jx_j$;
and $P(S)$ is a non-empty bounded polyhedron.
\end{lemma}

\begin{proof}
$x^{*}$ is integral with $x^{*}_j\ge1$ on $S$, so $x^{*}\in P(S)$. Conversely any
$x\in P(S)$, integral or not, has $x_j\ge1>0$ exactly on $S$, hence
$\Phi(x)=\sum_{j\in S}(s_j+c_jx_j)$; the first sum is constant on $P(S)$, so this is
affine there. Boundedness is immediate from $0\le x\le u$.
\end{proof}

The switch at the origin prevents $\Phi$ from being convex or concave as a whole.
Conditioning on the support freezes this switch: the fixed charges become a constant,
and nothing non-linear remains. No curvature is needed, so no sign
restriction is needed either.

\begin{theorem}\label{thm:main}
Suppose every non-empty support cell $P(S)$ is integral. Problem \eqref{eq:P} has an
optimal solution $\hat x$ that is a vertex of $P(\supp(\hat x))$, for fixed charges and
marginal rates of arbitrary sign.
\end{theorem}

\begin{proof}
Let $x^{*}$ be optimal and $S=\supp(x^{*})$. By Lemma~\ref{lem:cell}, $P(S)$ is a
non-empty bounded polyhedron on which $\Phi$ is affine, so $\Phi$ attains its minimum
over $P(S)$ at a vertex $\hat x$, and $\Phi(\hat x)\le\Phi(x^{*})$. Cell integrality
makes $\hat x$ integral, and Lemma~\ref{lem:cell} makes it feasible with support $S$, so
$\Phi(\hat x)\ge\OPT=\Phi(x^{*})$. Hence $\hat x$ is optimal and is a vertex of
$P(S)=P(\supp(\hat x))$.
\end{proof}

\begin{corollary}\label{cor:tu}
If $A$ is totally unimodular, $b$ is integral, and $u$ is integral, the hypothesis of
Theorem~\ref{thm:main} holds.
\end{corollary}

\begin{proof}
Orient greater-than rows by multiplying them by $-1$. Adjoining the signed identity
rows that define $x_j=0$ and $1\le x_j\le u_j$ preserves total unimodularity. Every
right-hand side is integral, so Hoffman--Kruskal makes each non-empty $P(S)$ integral.
\end{proof}

\begin{corollary}\label{cor:count}
Let $\hat x$ be as in Theorem~\ref{thm:main}, let $A_{\mathrm{t}}$ be the rows of $A$
tight at $\hat x$, and let $Q=\{j:1<\hat x_j<u_j\}$. Then the columns of
$A_{\mathrm{t}}$ indexed by $Q$ are linearly independent; in particular
$|Q|\le\operatorname{rank}A_{\mathrm{t}}$.
\end{corollary}

\begin{proof}
Suppose $d\ne0$ is supported on $Q$ with $A_{\mathrm{t}}d=0$. Strict interiority and
the positive slack of the finitely many inactive rows give $\varepsilon>0$ for which
$\hat x\pm\varepsilon d$ satisfy every row and bound; coordinates in $Q$ are strictly
interior; and coordinates outside $S$ are untouched, since $Q\subseteq S$. So $\hat x$
is the midpoint of a segment in $P(S)$, contradicting that it is a vertex.
\end{proof}

\begin{proposition}\label{prop:zangwill}
Let $X=[0,u_j]$ and let $\varphi_j$ be as in \eqref{eq:fc}. Then $\varphi_j$ is piecewise
concave in the sense of \cite{zangwill67} \emph{if and only if} $s_j\ge0$. His Theorem~6
therefore applies to a fixed charge that is paid and to no fixed charge that is received.
\end{proposition}

\begin{proof}
Zangwill calls $F$ \emph{continuous piecewise concave} (CPC) on $X$ when
$F=\max_i g_i$ for finitely many continuous concave $g_i$. Its \emph{basic sets} are
$B_i=\{z\in X: F(z)=g_i(z)\}$, and its \emph{terminal set} is
$T=\bigcup_{a\ne b}(B_a\cap B_b)$. The function $F$ is \emph{piecewise concave} (PC) when it is the pointwise limit of CPC functions
$\{F_k\}$ that share one terminal set $T$ and satisfy $F_k=F_1$ on $E[X]\cup T$.

Suppose $\{F_k\}$ is such a sequence for $\varphi_j$. Each $F_k$ is a maximum of finitely
many continuous functions, hence continuous; and $F_k=F_1$ on $E[X]\cup T$ together with
$F_k\to\varphi_j$ pointwise gives $F_1=\varphi_j$ there.

Since $0\in E[X]$ we get $F_1(0)=\varphi_j(0)=0$. As $F_1$ is continuous and
$s_j+c_jx\to s_j\ne0$ as $x\downarrow0$, choose $\delta\in(0,u_j)$ with $|F_1|<|s_j|/2$ and
$|s_j+c_jx|>|s_j|/2$ on $[0,\delta]$. Any $x\in T\cap(0,\delta)$ would satisfy
$F_1(x)=\varphi_j(x)=s_j+c_jx$, violating both bounds; hence $T\cap(0,\delta)=\emptyset$.

Fix $k$. Restricted to $(0,\delta)$, the basic sets form a finite, disjoint, relatively
closed cover. Connectedness therefore leaves exactly one non-empty member, so $F_k$
agrees there with a single continuous concave $g^k$. By continuity,
$g^k(0)=F_k(0)=F_1(0)=0$. Concavity with $g^k(0)=0$ gives, for $x\in(0,\delta)$ and
$\lambda\in(0,1)$,
\[
g^k(\lambda x)\ \ge\ \lambda g^k(x)+(1-\lambda)g^k(0)\ =\ \lambda g^k(x).
\]
Letting $k\to\infty$ and using $g^k\to\varphi_j$ pointwise on $(0,\delta)$,
\[
s_j+c_j\lambda x\ \ge\ \lambda(s_j+c_jx),\qquad\text{that is}\qquad s_j(1-\lambda)\ \ge\ 0 ,
\]
so $s_j\ge0$, contradicting $s_j<0$.

Conversely let $s_j\ge0$, put $K_0=s_j/u_j+|c_j|+1$ and
\[
F_k(x)\ =\ \min\bigl\{(k+K_0)x,\ s_j+c_jx\bigr\},\qquad k\ge1 .
\]
Each $F_k$ is a minimum of two affine functions, hence continuous and concave, hence CPC as
a family with the single member $F_k$ itself. That family has no second member, so no two
basic sets meet and the terminal set is $\emptyset$ for every $k$; condition~(2) then only
asks for agreement on $E[X]=\{0,u_j\}$. There $F_k(0)=\min\{0,s_j\}=0=\varphi_j(0)$, using
$s_j\ge0$, and $(k+K_0)u_j>s_j+c_ju_j$ for every $k\ge1$ by the choice of $K_0$, so
$F_k(u_j)=s_j+c_ju_j$ independently of $k$. For fixed $x\in(0,u_j]$, the inequality
$(k+K_0)x>s_j+c_jx$ holds once $k$ is large. Hence $F_k(x)\to s_j+c_jx$, while
$F_k(0)=0$. Thus $F_k\to\varphi_j$ pointwise, and $\{F_k\}$ is a defining sequence.
\end{proof}

The sign of the fixed charge is the boundary of the class covered by Zangwill's theory.
The classical route applies to a paid charge and closes when the charge is received.

\section{What does not survive}\label{sec:fails}

Losing a hypothesis is not losing a conclusion. A received charge is not piecewise
concave, so \cite{fk71,love73} and the recursions built on them no longer \emph{apply};
that alone would leave the classical properties unproven rather than false. What follows
is stronger: in each setting the conclusion itself is violated.

Corollary~\ref{cor:count} constrains only the coordinates with $1<\hat x_j<u_j$. The
classical structural properties constrain those with $0<\hat x_j<u_j$. Under negative
charges this difference is decisive: opening is rewarded and positive volumes can be
driven to their minimum level. Each witness below has a
\emph{unique} optimum, so no choice of optimal solution repairs the classical statement,
and each respects the coupling $s_j=p_j\sigma_j$, $c_j=p_j\gamma_j$ with
$\sigma_j,\gamma_j>0$ of Section~1. Thus the failures do not arise from letting the
fixed charge and the marginal rate carry opposite signs.

\paragraph{Lot sizing.} Let $x_t$ be production in period $t$, $u_t=C$, and let the rows
of $A$ be the prefix sums bounded below by cumulative demand and above by cumulative
demand plus a storage bound; $A$ is an interval matrix, hence totally unimodular. Take
$T=3$, $C=3$, zero demand, storage bound $2$, and prices $p=(-1,-1,1)$ with
$\sigma_t=\gamma_t=1$, so $s_t=c_t=p_t$. The unique optimum is $x=(1,1,0)$, with prefix
sums $(1,2,2)$ and value $-4$. Rows $2$ and $3$ are tight and row $1$ is not, so no
inventory point separates periods $1$ and $2$, both of which produce strictly between $0$
and $C$: the fractional-period property of \cite{fk71,love73} fails on the unique
optimum. Here $Q=\emptyset$, so Corollary~\ref{cor:count} holds vacuously.

\paragraph{Fixed-charge transportation.} Let $A$ be the $2\times2$ bipartite incidence
matrix with supplies and demands $(2,2)$, $u_{ij}=2$ and $s_{ij}=c_{ij}=-1$. The unique
optimum, of value $-8$, sends one unit on each of the four arcs. No arc is strictly interior, so
again Corollary~\ref{cor:count} is vacuous, while the four positive arcs exceed the
classical bound $m+n-1=3$ and form a cycle.

\paragraph{Fixed-charge min-cost flow.} Read the same instance as a four-node flow
problem. Its four positive arcs exceed the spanning-tree bound $n-1=3$.

So the sign hypothesis cannot be removed from the classical statements; it can only be
removed from the vertex statement, which is weaker. Section~\ref{sec:alg} shows the
difference is not academic.

\section{A classical state space that does not transfer}\label{sec:alg}

Corollary~\ref{cor:count} might suggest that an algorithm based only on a vertex argument
transfers to arbitrary signs. The recursion of \cite{apr15} gives a concrete test: it uses
the stronger classical property rather than the surviving vertex property.

The recursion of \cite{apr15} is $O(T^4)$ and is correct under the concave costs it
assumes. Its state space does not, however, transfer unchanged to arbitrary signs. It is
built directly on Love's property, which it states, after
\cite{love73}, in terms of periods that are \emph{fractional} in the classical sense
$0<x_t<P$: within a subplan, every period produces either zero or $P$, apart from at most
one fractional period. Take $T=4$, $C=5$, storage
bound $4$, demand $(0,0,0,4)$, zero initial and terminal inventory, and $s_t=c_t=-1$.
Conservation fixes total production at four, so every feasible plan costs
$-4-|\supp(x)|$. The unique optimum is $(1,1,1,1)$, of value $-8$. Before final demand,
inventory cannot return to zero after positive production; reaching the upper bound four
also consumes all production. A plan decomposable into the subplans of \cite{apr15}
therefore has at most one positive period, necessarily of size four, and costs $-5$.
Thus the state space cannot express the optimum outside its concavity assumptions; the
absolute gap is three cost units.

\section{Support-cell integrality cannot be discarded}\label{sec:sharp}

\begin{proposition}\label{prop:sharp}
There is an instance of \eqref{eq:P} with a non-integral support cell whose unique optimum
is not a vertex of that cell, so the conclusion of Theorem~\ref{thm:main} fails.
\end{proposition}

\begin{proof}
Take the odd-cycle system $x_1+x_2\le6$, $x_2+x_3\le7$, $x_1+x_3\le6$, whose constraint
matrix has determinant $2$, with $u=(5,5,5)$, $s=(1.32,\,2.51,\,0.55)$ and
$c=(-2.02,\,-2.06,\,-1.19)$. Exhaustive enumeration over the $113$ feasible integer
points gives the unique optimum $\hat x=(2,4,3)$ of value $-11.47$, the runner-up being
$(3,3,3)$ at $-11.43$. Its support is $\{1,2,3\}$, so its cell is
$\{x:\ Ax\le b,\ 1\le x\le5\}$. The direction $d=(1,-1,1)$ preserves the first two
rows, and $\hat x\pm d/2$ satisfy the third row and every cell bound. These two distinct
cell points have midpoint $\hat x$, so $\hat x$ is not a vertex.
\end{proof}

This shows that support-cell integrality cannot be removed without replacement; it does
not make total unimodularity necessary.

\begin{remark}[A modelling consequence]\label{rem:model}
The natural mixed-integer formulation of \eqref{eq:fc} writes $x_j\le u_jy_j$ with $y_j$
binary. For $s_j\ge0$ that suffices, because setting $y_j=1$ with $x_j=0$ only adds
cost. For $s_j<0$ it does not: the solver sets $y_j=1$ and $x_j=0$ and banks the charge.
Over integer $x$ the constraint $y_j\le x_j$ is required. Consider four periods with
prices $(1,-5,-5,-5)$, two units due by the horizon, $u=2$, and unit coefficients. The
guarded formulation returns $-20$, whereas the unguarded one returns $-25$. The latter
creates an artificial credit of $5$ by ``using'' idle activities to collect the charge.
The object it optimises is the lower
semicontinuous hull of $\varphi_j$, which is the function \cite{kya14} assume and
\eqref{eq:fc} is not.

The unguarded pattern is the standard one. \cite{apr15} write $0\le x_t\le Py_t$ with no
lower link, and \cite{lcw01} write $x_t\le My_t$ having just \emph{defined}
``$y_t=1$ if $x_t>0$, and $0$ otherwise'', a two-sided definition implemented by a
one-sided constraint. Both are correct in their own settings, where the setup cost is
non-negative and $y_t=1$ with $x_t=0$ is never optimal; we report no error in either.

Two qualifications keep this in proportion. First, energy models already use the repair.
The minimum-load constraint
$p_t\ge P^{\min}z^{\mathrm{on}}_t+P^{\mathrm{sb}}z^{\mathrm{sb}}_t$ in
\cite{raheli23} forbids an activated unit from consuming nothing. What is missing
from the lot-sizing formulations above is a constraint the energy literature writes as a
matter of course. Second, the general shape is widespread and not confined to
\eqref{eq:fc}. The same paper \cite{raheli23}, whose numerical study is explicitly about
nonpositive prices, defines its start-up binary by
$z^{\mathrm{su}}_t\ge z^{\mathrm{off}}_{t-1}+z^{\mathrm{on}}_t+z^{\mathrm{sb}}_t-1$ with no
upper pin, which is correct only because the start-up cost is a positive constant rather
than a price. Our point is about \eqref{eq:fc}; the pattern is broader.
\end{remark}

\section{Computational verification}

The supplement reproduces every finite witness by exhaustive enumeration. This includes
the instance in Section~\ref{sec:alg}. It also checks the vertex theorem on
$1{,}085$ independently generated matrices verified to be totally unimodular. A negative
control allows supports to vary and breaks the affine identity on $217$ of $220$ instances.
These computations test the implementations and transcriptions; the general claims follow
from the proofs above.

The supplement also contains a Lean 4 formalization of the support-conditioning identity,
the TU support-cell integrality argument, the vertex theorem, and the active-column
corollary. The TU proof selects a nonsingular active minor and derives integrality from its
unit determinant. The formal theorem takes an integer TU constraint matrix, integral
right-hand sides and upper bounds, and feasibility as inputs. It produces an optimal
integer solution that is a vertex of its support cell. The general theorem chain contains
no native evaluation or project-added axioms.

\section{Concluding remarks}

The opening question has a narrow answer. A vertex optimum survives arbitrary signs after
conditioning on an integral support cell. The stronger classical structures do not.
Zangwill's class ends at the same sign change.

The vertex result follows from known integrality. It continues the vertex viewpoint
of Hirsch and Hoffman \cite{hh61}, Hirsch and Dantzig \cite{hd68}, and Murty
\cite{murty68}. It also relates to results for nonconvex production costs
\cite{jr73,sr82} and implied integrality \cite{hw26}. We claim no new source of
integrality.

Signed setup coefficients also appear in Van Hoesel and Wagelmans:
\cite{vhw93} allow a negative setup cost and note it is then profitable to set up even
with no production, absorbing such terms as constants. Their formulation lets a
setup be chosen while $x=0$, which is a different object from the charge-if-positive
function studied here.

The stronger classical statements fail in all three settings. A
published polynomial recursion, correct under its own hypotheses, loses optimality when
its state space is transferred unchanged to a four-period signed-charge instance. The
practical conclusion is direct. A negative fixed charge changes which plans are optimal:
few-and-large can become many-and-minimal.

The intrinsic condition in Theorem~\ref{thm:main} is integrality of the support cells;
characterising broader matrix classes that guarantee it is left open. Of the formulations
we checked, the two lot-sizing ones \cite{apr15,lcw01} omit the
guard of Remark~\ref{rem:model}. This omission is harmless under their assumptions but
unsafe under signed charges. The energy model \cite{raheli23} carries its continuous analogue as a
minimum-load constraint. The exposure we describe is therefore a property of how these models would be
combined, rather than one we exhibit in print.

\begin{thebibliography}{99}
\bibitem{ak05} A. Atamt\"urk, S. K\"u\c{c}\"ukyavuz, Lot sizing with inventory bounds and
fixed costs: polyhedral study and computation, \emph{Operations Research} 53 (4) (2005)
711--730.
\bibitem{apr15} A. Akbalik, B. Penz, C. Rapine, Capacitated lot sizing problems with
inventory bounds, \emph{Annals of Operations Research} 229 (1) (2015) 1--18.
\bibitem{fk71} M. Florian, M. Klein, Deterministic production planning with concave costs
and capacity constraints, \emph{Management Science} 18 (1) (1971) 12--20.
\bibitem{flr80} M. Florian, J. K. Lenstra, A. H. G. Rinnooy Kan, Deterministic production
planning: algorithms and complexity, \emph{Management Science} 26 (7) (1980) 669--679.
\bibitem{hd68} W. M. Hirsch, G. B. Dantzig, The fixed charge problem, \emph{Naval
Research Logistics Quarterly} 15 (3) (1968) 413--424.
\bibitem{hh61} W. M. Hirsch, A. J. Hoffman, Extreme varieties, concave functions, and
the fixed charge problem, \emph{Communications on Pure and Applied Mathematics} 14 (3)
(1961) 355--369.
\bibitem{hw26} R. van der Hulst, M. Walter, Implied integrality in mixed-integer
optimization, \emph{Mathematical Programming} (2026), in press.
\bibitem{jr73} R. Jagannathan, M. R. Rao, A class of deterministic production planning
problems, \emph{Management Science} 19 (11) (1973) 1295--1300.
\bibitem{kya14} E. Koca, H. Yaman, M. S. Akt\"urk, Lot sizing with piecewise concave
production costs, \emph{INFORMS Journal on Computing} 26 (4) (2014) 767--779.
\bibitem{lcw01} C.-Y. Lee, S. \c{C}etinkaya, A. P. M. Wagelmans, A dynamic lot-sizing
model with demand time windows, \emph{Management Science} 47 (10) (2001) 1384--1395.
\bibitem{love73} S. F. Love, Bounded production and inventory models with piecewise
concave costs, \emph{Management Science} 20 (3) (1973) 313--318.
\bibitem{murty68} K. G. Murty, Solving the fixed charge problem by ranking the extreme
points, \emph{Operations Research} 16 (2) (1968) 268--279.
\bibitem{ou17} J. Ou, Improved exact algorithms to economic lot-sizing with piecewise
linear production costs, \emph{European Journal of Operational Research} 256 (3) (2017)
777--784.
\bibitem{raheli23} E. Raheli, Y. Werner, J. Kazempour, A conic model for electrolyzer
scheduling, \emph{Computers \& Chemical Engineering} 179 (2023) 108450.
\bibitem{seel21} J. Seel, D. Millstein, A. Mills, M. Bolinger, R. Wiser, Plentiful
electricity turns wholesale prices negative, \emph{Advances in Applied Energy} 4 (2021)
100073.
\bibitem{sr82} S. Sadagopan, A. Ravindran, A vertex ranking algorithm for the
fixed-charge transportation problem, \emph{Journal of Optimization Theory and
Applications} 37 (2) (1982) 221--230.
\bibitem{swoveland75} C. Swoveland, A deterministic multi-period production planning
model with piecewise concave production and holding-backorder costs, \emph{Management
Science} 21 (9) (1975) 1007--1013.
\bibitem{vhw93} C. P. M. Van Hoesel, A. P. M. Wagelmans, Sensitivity analysis of the
economic lot-sizing problem, \emph{Discrete Applied Mathematics} 45 (3) (1993) 291--312.
\bibitem{zangwill67} W. I. Zangwill, The piecewise concave function, \emph{Management
Science} 13 (11) (1967) 900--912.
\end{thebibliography}
\end{document}
